\section{Probability Space}

\subsection{Random Experiment}

\begin{itemize}
\item A \textit{\textbf{random experiment}} is an experiment in which the outcome varies in an unpredictable fashion when the experiment is repeated under the same conditions, i.e., \textit{random}.
\item An \textit{\textbf{outcome}} of a random experiment is a result that cannot be decomposed into other results.
\end{itemize}

\subsection{Sample Space}

\begin{itemize}
\item The set of all possible outcomes for a random experiment is called the \textit{\textbf{sample space}} $\Omega$.
\end{itemize}

\begin{example}\begin{enumerate}
\item Tossing a coin: $\Omega = \{H,T\}$
\item Tossing a coin twice: $\Omega = \{HH,HT,TH,TT\}$
\item Taking EEE 5544 and getting a grade:
$\Omega = \{A,A -,B+,B,B -,C+,C,D,F\}$
\item Picking an integer at random:
$\Omega = \{\ldots, -3, -2, -1,0,1,2,3,\ldots\}$
\item Checking the temperature (in $K$) at a random location:
$\Omega = (0,\infty)$
\end{enumerate}

\begin{itemize}
\item \textbf{Finite} $\Omega$: Examples 1, 2, and 3
\item \textbf{Countably infinite} $\Omega$: Example 4
\item \textbf{Discrete} $\Omega$: Examples 1--4
\item \textbf{Continuous (uncountably infinite)} $\Omega$: Example 5
\end{itemize}

\end{example}\subsection{Event Class}

\begin{itemize}
\item An \textit{\textbf{event}} is a set of outcomes, i.e., a subset of $\Omega$,
whose ``likelihood'' can be quantified.
\item Two special events:

\begin{itemize}
\item $\Omega =$ \textit{\textbf{certain event}}
\item $\emptyset =$ \textit{\textbf{impossible event}}
\end{itemize}


\item An \textit{\textbf{event class}} $\mathcal{F}$ of a random experiment is a collection of events.
\end{itemize}

\subsubsection{Discrete Sample Space}

For a discrete $\Omega$, we typically consider the event class
consisting of all subsets of $\Omega$, i.e.,
$\mathcal{F} = 2^\Omega.$
This is called the \textit{power set} of $\Omega$.

\begin{example}\begin{enumerate}
\item Tossing a coin:
$\mathcal{F} = \{\emptyset,\{H\},\{T\},\{H,T\}\}.$
\item Tossing a coin twice:
\end{enumerate}
\begin{equation*}
\mathcal{F}
   =
   \left\{
   \emptyset,
   \{HH\},
   \{TT\},
   \{HT\},
   \{TH\},
   \{HH,TT\},
   \{HH,HT\},
   \ldots,
   \{HH,TT,HT\},
   \ldots,
   \{HH,TT,HT,TH\}
   \right\}.
\end{equation*}

\begin{enumerate}[resume]
\item If $\Omega=\{a_1,a_2,\ldots,a_N\}$, then

\begin{equation*}
\mathcal{F}
   =
   \left\{
   \emptyset,
   \{a_1\},
   \ldots,
   \{a_N\},
   \{a_1,a_2\},
   \ldots,
   \{a_1,a_2,a_3\},
   \ldots,
   \{a_1,a_2,\ldots,a_N\}
   \right\}.
\end{equation*}

Altogether, there are $2^N$ events (why?).


\item For a countably infinite sample space
$\Omega=\{a_1,a_2,\ldots\},$
we can use the same idea as in the finite case, using the power set as the event space:
\begin{equation*}
\mathcal{F}
   =
   2^\Omega
   =
   \left\{
   \emptyset,
   \{a_1\},
   \ldots,
   \{a_N\},
   \{a_1,a_2\},
   \ldots,
   \{a_1,a_2,a_3\},
   \ldots,
   \{a_1,a_2,\ldots,a_n\},
   \ldots
   \right\}.
\end{equation*}

This enumerates all subsets of $\Omega$. However, we need a better way to describe the power set.
\end{enumerate}

\end{example}\subsubsection{Continuous Sample Space}

\begin{itemize}
\item For a continuous sample space, such as $\Omega=(0,\infty)$, there are uncountably infinitely many outcomes. Therefore, the enumeration method used in the finite case cannot be used to construct the event class.
\item To discuss the solution, we first need to study set operations.
\end{itemize}

\subsection{Set Operations}

\begin{itemize}
\item \textbf{Union}:
$A\cup B=\{x\in\Omega:x\in A\text{ or }x\in B\}$
\item \textbf{Intersection}:
$A\cap B=\{x\in\Omega:x\in A\text{ and }x\in B\}$
\item \textbf{Complement}:
$A^c=\{x\in\Omega:x\notin A\}$
\item \textbf{Difference}:
$A \setminus B =
\{x\in\Omega:x\in A \text{ and } x\notin B\}
= A\cap B^c$
\end{itemize}

\begin{remark}[De Morgan’s Laws]\begin{enumerate}
\item $(A\cup B)^c=A^c\cap B^c$.
\item $(A\cap B)^c=A^c\cup B^c$.
\end{enumerate}

\end{remark}\begin{proof}We prove the first identity as an example.

Since
\begin{align*}
   x\in(A\cup B)^c
   &\Rightarrow x\notin A\cup B \\
   &\Rightarrow x\notin A\text{ and }x\notin B \\
   &\Rightarrow x\in A^c\text{ and }x\in B^c \\
   &\Rightarrow x\in A^c\cap B^c,
   \end{align*}
we have $(A\cup B)^c\subseteq A^c\cap B^c$.

Conversely, since
\begin{align*}
   x\in A^c\cap B^c
   &\Rightarrow x\in A^c\text{ and }x\in B^c \\
   &\Rightarrow x\notin A\text{ and }x\notin B \\
   &\Rightarrow x\notin A\cup B \\
   &\Rightarrow x\in(A\cup B)^c,
   \end{align*}
we have $A^c\cap B^c\subseteq(A\cup B)^c$.

Combining the two inclusions gives
$(A\cup B)^c=A^c\cap B^c$.

\end{proof}\begin{remark}Union and complement are enough to specify all other set operations.

\end{remark}\subsection{Field (Algebra)}

\begin{itemize}
\item A \textbf{field (algebra)} is a collection of subsets of $\Omega$ that is closed under union and complement and contains both $\emptyset$ and $\Omega$. Thus, $\mathcal{F}$ is a field if it satisfies the following conditions:

\begin{enumerate}
\item $\emptyset\in\mathcal{F}$ and $\Omega\in\mathcal{F}$.
\item If $A\in\mathcal{F}$ and $B\in\mathcal{F}$, then $A\cup B\in\mathcal{F}$.
\item If $A\in\mathcal{F}$, then $A^c\in\mathcal{F}$.
\end{enumerate}


\item Operationally, a field contains all subsets of $\Omega$ that can be obtained by applying a finite number of set operations.
\end{itemize}

\begin{remark}For any finite $\Omega$, the power set $2^\Omega$ is a field.

\end{remark}\subsection{$\sigma$-Field ($\sigma$-Algebra)}

\begin{itemize}
\item A field $\mathcal{F}$ is a \textbf{$\sigma$-field} if it satisfies the additional condition

\begin{enumerate}[resume]
\item If $A_1,A_2,\ldots\in\mathcal{F}$, then
$\bigcup_{i=1}^{\infty}A_i\in\mathcal{F}$.
\end{enumerate}


\item Operationally, a $\sigma$-algebra contains all subsets of $\Omega$ that can be obtained through a countable number of set operations.
\end{itemize}

\begin{remark}For countably infinite $\Omega$, $2^\Omega$ is a $\sigma$-field.

\end{remark}\subsubsection{Generating a $\sigma$-Field}

\begin{itemize}
\item For a continuous (uncountably infinite) sample space $\Omega$, we often \textit{generate} $\sigma$-fields from collections of subsets of $\Omega$.
\item For a collection $\mathcal{A}$ of subsets of $\Omega$, consider the $\sigma$-field generated by $\mathcal{A}$: the smallest $\sigma$-field containing every subset that can be formed by performing countable set operations on members of $\mathcal{A}$.
\item The $\sigma$-field generated by $\mathcal{A}$ is denoted by $\sigma(\mathcal{A})$.
\end{itemize}

\subsection{Event Class for a Continuous Sample Space}

\begin{itemize}
\item For a continuous (uncountably infinite) sample space $\Omega$, we consider event classes that are $\sigma$-fields on $\Omega$.
\end{itemize}

\begin{example}[Event class for ]\textbf{$\Omega=\mathbb{R}$}\\
\begin{itemize}
\item Consider the collection $\mathcal{A}$ of all subsets of $\mathbb{R}$ of the form $( -\infty,a)$ for some $a\in\mathbb{R}$.
\item The $\sigma$-field generated by $\mathcal{A}$ is called the \textit{Borel $\sigma$-field}. It is commonly used as an event class for the sample space $\mathbb{R}$.
\item The Borel $\sigma$-field contains essentially all events of practical interest.
\item Although the \textit{Borel sets} form a very rich class of subsets of $\mathbb{R}$, some subsets of $\mathbb{R}$ are not Borel sets. This is critically important for allowing the likelihoods of events to be quantified.
\end{itemize}

\end{example}\subsection{Description of a Random Experiment}

\begin{itemize}
\item The sample space $\Omega$ and an event class $\mathcal{F}$ can be used to describe a random experiment mathematically.
\item We also need a way to quantify the likelihood of events. This is done by a mapping called a \textit{probability measure} $P$, which will be discussed in the next section.
\item In summary, the mathematical object describing a random experiment is the triple
$(\Omega,\mathcal{F},P)$,
called the \textit{\textbf{probability space}} of the random experiment.
\end{itemize}